\documentclass[11pt]{article}
\usepackage{fontspec}
\setmainfont{Times New Roman}
\setsansfont{Arial}
\setmonofont{Consolas}
\usepackage{amsmath,amssymb}
\usepackage{geometry}
\usepackage{microtype}
\geometry{a4paper,margin=3cm}
\setlength{\parindent}{1.25em}
\setlength{\parskip}{0pt}
\emergencystretch=4em

\begin{document}

\begin{center}
\textbf{§ 3. \quad Redukcijne teoremy.}
\end{center}

Znane teoremy o redukciji form i sistemov form trěba nyně svęzati s paragrafami 1 i 2; za to potrěbne sut něktore opreděljenja, vzęte iz teorije binarnyh form.

a) Kako {\itshape relativno polny sistem} po modulu danogo ręda form označajemo sistem form s takym svojstvom: vse tvorjenja, ktore nastavajut črez svijanje libovolnogo produkta tyh form, možno izraziti kako cěle racionalne funkcije form sistema, do dodatnogo člena izrazov, ktore nastali črez svijanje sistemnyh form so sistemom modula\footnote{Gordan--Kerschensteiner, \emph{Vorlesungen über Invariantentheorie}, tom II, s. 227.}.

b) Formu nazyvajemo {\itshape reducibilnoju}, kogda ju možno izraziti črez formy, ktore imajut invarianty kako faktory, ili črez «vyše formy»; to jest črez vyše formy cělogo ręda form, k ktoromu ona naleži, ili črez formy, ktore soderžat simboly modula vo vyšem porędku.

Teoremy o reducentah\footnote{Gordan, \emph{Math. Annalen}, tom XVII, s. 222.} možno v najobčejšoj formě izrekti tak:

Opreděljenje: {\itshape reducent jest reducibilny ręd form.}

{\itshape Teorema II. Jestli početna forma ręda form jest reducibilna zato, že jeden iz jej členov nastal črez svijanje s reducentom (ima reducent kako faktor), i jestli konečna forma ręda form nastala iz togo samogo člena črez svijanje, togda cěly ręd form jest reducibilny}\footnote{Uproščenje, ktore nastaje črez teoremu II, možno pokazati na slědujućem primjeru. Za specialnu formu $f=x_1^3x_2+x_2^3x_3+x_3^3x_1$ izraz $a_y^2a_x^2u_y^2$ jest reducent. Iz togo po teoremě II slěduje redukcija ręda form
\[
\theta_\nu(\vartheta\nu x)\theta_x^3u_\vartheta u_\nu^2
=
a_\nu^2(abu)-a_\nu b_\nu(aba),
\]
bo $\theta_\nu^4=a_\nu^2b_\nu^2(abu)^2$ nastala iz člena $a_\nu^2(abu)$ črez svijanje. U Gordana, tamže, s. 231, nasuproti tomu, redukcija form
\[
\theta_\nu(\vartheta\nu x),\quad \theta_\nu(\vartheta\nu x)^2,\quad
\theta_\nu^2,\quad \theta_\nu^2(\vartheta\nu x),\quad
\theta_\nu^2(\vartheta\nu x)^2
\]
provodit se pojedinačno.}.

Dokaz: Ręd form nastaje iz početnogo člena črez svijanje samogo sebe, to jest s jednej strany črez vyše svijanja s reducentom, s drugoj strany črez svijanje reducenta samogo sebe. V obojih slučajah, po § 2, možno vse formy ręda form predstaviti kako polary (prěkrivanja) nad rędom form reducenta.

Zaključenje od početnoj formy k cělomu rędu form uže ne plati pri ostalyh metodah redukcije. To sut:

1) tak nazyvana «dvojna redukcija».

Pod tym razumějemo redukciju izraza, ktory ima dva medžusobno nezavisne reducenty kako faktory, dvojim sposobom; črez to nastavajut relacije medžu vyšimi formami, k ktorym bylo reducirano.

Sistematičnym priměnjenjem «dvojnoj redukcije» trěba s jednej strany dobiti vse relacije\footnote{Tako sut, na primjer, črez «dvojnu redukciju» najdene: redukcija formy $a_\eta^4$ (§ 10), jedinoj formy, ktoru Maisano zbytočno vnesl v sistem; potom v primětki k uvodu spomenute relacije medžu 3 kovariantami i 3 invariantami.}; s drugoj strany, jestli «dvojna redukcija», priměnjena k različnym izrazam, ne daje novih relacij, imamo i kriterij nezavisnosti vyših form i kontrolu računa.

2) {\itshape Svijanje s razložimymi formami.}

Pod «razložimymi formami» -- s malym občenjem običnogo pojma -- razumějemo formy, ktore možno izraziti črez produkty form nižšego stepena i črez «vyše» formy. Produkty form pri jednokratnom svijanju prěhodet v produkty; takože produkty nelinearnyh form pri dvokratnom svijanju pri specializacijah a') i b') (§ 2) po teoremě produkta:
\[
a_yb_y a_xb_x=\frac12a_y^2b_x^2+\frac12b_y^2a_x^2-\frac12(ab\widehat{yx})^2.
\]
Prve polary (prěkrivanja) i něktore druge polary razložimyh form zato sut znovu razložime formy. Iz togo slěduje:

Člen jednokratnogo (odnošno dvokratnogo) prěkrivanja nad razložimoju formoju, ktory nastal črez jednokratno (odnošno dvokratno) svijanje s razložimoju formoju, sam staje razložima forma:

a) kogda slědujuće vyše formy v rędu form razložimoj formy razlagajut se ili stavaju reducibilne, ili takože kogda pri jednokratnom svijanju nastavajut specializacije $y=\widehat{uv}$ ili $v=\widehat{xy}$;

b) kogda ostale jednokratne svijanja vode k vyšim formam (srav. § 12, B. $H_k$ i $H_k(k\eta x)$).

Takovy člen prěkrivanja može takože razložiti se:

c) kogda ostale jednokratne svijanja vode k nižšim formam, ktore nezavisno od svijanja s razložimoju formoju razlagajut se, to jest možno ih izraziti črez produkty i črez formu, ktoru trěba reducirati. Pri tom trěba vsakokratno najprvo računom ověrjati, že čislovy koeficient odpovědajučego člena ne isčezaje identično. To jest ničto drugo než «dvojna redukcija» nižšej formy (srav. § 23, C. $H_q(\theta H u)(\theta s u)$).

Razložime formy nastavaju črez vse jednokratne svijanja s funkcionalnymi determinantami\footnote{Srav. Gordan, tamže, § 3.}, kromě togo črez něktore vyše svijanja. Imamo
\[
(st\widehat{yx})s_y=\frac12\,t\,s_y^2-\frac12\,s\,t_y^2+\frac12(st\widehat{yx})^2,
\]
\[
(st\widehat{yx})^2s_y^2=\frac13\,t\,s_y^3-\frac13\,s\,t_y^3+s_y(st\widehat{yx})^2-\frac13(st\widehat{yx})^3.
\]
Analogične formuly platet za dualistične formy
\[
(\sigma\tau\widehat{\nu u})v_\sigma
\quad\text{i}\quad
(\sigma\tau\widehat{\nu u})v_\sigma^2.
\]

Iz teoremov togo paragrafa za postavjenje vsih nereducibilnyh form, ktore nastavaju iz svijanja $S$ s $T$, dobiva se slědujuće pravilo:

Iděmo, počinajući s najnižšimi svijanjami, po libovolnom naprijed ustanovjenom putu (na primjer po rędah ili simetrično k diagonalnomu členu) tako daleko v tvorjenju ręda form $ST$, poky ne dojdemo do formy, ktora ima reducent kako faktor. Ręd form opreděljeny toju formoju trěba izpustiti, kogda konečna forma izpolnja uslovje postavjeno v teoremě II. Poslě izključenja vsih reducibilnyh rędov form trěba črez priměnjenje dvojnoj redukcije odstraniti vse ostale reducibilne, i takože vse razložime formy. Města cělogo ręda form, ktore sut dvojno prěkryte nezavisno reducibilnymi rędami form, davajut povod k redukciji vyših form (srav. § 11, D).

\clearpage
\noindent\textbf{Teorema III.}
{\itshape Formy vzęte iz binarnoj oblasti, to jest one formy, ktore nastavaju iz sistemnyh form odpovědajučego binarnogo formy po Clebschovom principu prěnosa črez obrubljenje, tvoręt relativno polny sistem mod $(abc)$.}

\noindent\textbf{Dokaz:}
Binarnoj relaciji, ktora govori, že formy $Q'_1,\ldots,Q'_n$ tvoręt polny sistem jednoj binarnoj osnovnoj formy:
\[
Q'-F(Q')=0=\sum_\lambda\bigl\{(ab)c_x+(bc)a_x+(ca)b_x\bigr\}\varphi_\lambda^{*}
\]
črez obrubljenje odgovarja ternarna relacija:
\[
Q-F(Q_i)=\sum_\lambda u_x\cdot(abc)\varphi_\lambda .
\]
To jest ale opreděljajuće ravnanje za relativno polny sistem mod $(abc)$. Za biquadratičnu formu sistem vzętih form sostaji iz slědujućih form:
\[
\begin{gathered}
f=a_x^4,\qquad
\theta=\theta_x^4u_\vartheta^2=(abu)^2a_x^2b_x^2,\qquad
K=K_x^6u_k^3=(a\theta u)u_\vartheta^2a_x^3\theta_x^3,\\
j=u_j^6=(a\theta u)^4u_\vartheta^2,\qquad
\nu=u_\nu^4=(abu)^4,
\end{gathered}
\]
medžu ktorymi prve četyri tvoręt relativno polny sistem mod $((abc),\nu)$.

\end{document}
